\documentclass[12pt,a4paper]{article}

% ── Packages ──────────────────────────────────────────────────────────────────
\usepackage[a4paper, margin=1in]{geometry}
\usepackage{graphicx}
\usepackage{booktabs}
\usepackage{array}
\usepackage{longtable}
\usepackage{tabularx}
\usepackage{xcolor}
\usepackage{colortbl}
\usepackage{fancyhdr}
\usepackage{titlesec}
\usepackage{enumitem}
\usepackage{hyperref}
\usepackage[T1]{fontenc}
\usepackage[utf8]{inputenc}
\usepackage{parskip}
\usepackage{amsmath}
\setlength{\headheight}{14pt}
\usepackage{multirow}
\usepackage{makecell}
\usepackage{lastpage}

% ── Colors ────────────────────────────────────────────────────────────────────
\definecolor{darkblue}{RGB}{31,56,100}
\definecolor{medblue}{RGB}{46,80,144}
\definecolor{lightblue}{RGB}{238,242,248}
\definecolor{tableheadbg}{RGB}{31,56,100}
\definecolor{rowaltbg}{RGB}{238,242,248}
\definecolor{hred}{RGB}{192,0,0}
\definecolor{mred}{RGB}{127,127,0}
\definecolor{lgreen}{RGB}{0,128,0}
\definecolor{gray}{RGB}{136,136,136}

% ── Hyperref Setup ────────────────────────────────────────────────────────────
\hypersetup{
    colorlinks=true,
    linkcolor=medblue,
    urlcolor=medblue,
    pdftitle={SRS - Hotel Automation Software v1.0},
    pdfauthor={Project Team}
}

% ── Section Title Formatting ──────────────────────────────────────────────────
\titleformat{\section}
  {\Large\bfseries\color{darkblue}}
  {\thesection}{1em}{}
  [\vspace{2pt}\color{medblue}\titlerule[1pt]\vspace{4pt}]

\titleformat{\subsection}
  {\large\bfseries\color{medblue}}
  {\thesubsection}{1em}{}

\titleformat{\subsubsection}
  {\normalsize\bfseries\color{medblue}}
  {\thesubsubsection}{1em}{}

% ── Header & Footer ───────────────────────────────────────────────────────────
\pagestyle{fancy}
\fancyhf{}
\fancyhead[L]{\small\color{medblue} Software Requirements Specification}
\fancyhead[R]{\small\color{gray} Hotel Automation Software}
\fancyfoot[R]{\small\color{gray} Page \thepage\ of \pageref{LastPage}}
\renewcommand{\headrulewidth}{0.5pt}
\renewcommand{\footrulewidth}{0.5pt}

% ── Table Column Types ────────────────────────────────────────────────────────
\newcolumntype{L}[1]{>{\raggedright\arraybackslash}p{#1}}
\newcolumntype{C}[1]{>{\centering\arraybackslash}p{#1}}
\newcolumntype{R}[1]{>{\raggedleft\arraybackslash}p{#1}}

% ── Table Header Macro ────────────────────────────────────────────────────────
\renewcommand{\thead}[1]{%
  \cellcolor{tableheadbg}\color{white}\textbf{#1}%
}

% ── Requirement Box Macro ─────────────────────────────────────────────────────
\newcommand{\req}[2]{%
  \noindent\textbf{\color{medblue}#1:} #2\par\smallskip
}

% ══════════════════════════════════════════════════════════════════════════════
\begin{document}

% ── COVER PAGE ────────────────────────────────────────────────────────────────
\thispagestyle{empty}
\begin{center}
  \vspace*{3cm}
  {\Huge\bfseries\color{darkblue} SOFTWARE REQUIREMENTS\\[0.4em] SPECIFICATION}\\[1em]
  {\LARGE\color{medblue} Hotel Automation Software}\\[3cm]

  \begin{tabular}{|L{4cm}|L{8cm}|}
    \hline
    \rowcolor{tableheadbg}
    \color{white}\textbf{Document Type} & \color{white} Software Requirements Specification (SRS) \\
    \hline
    \rowcolor{rowaltbg}
    \textbf{Subject}           & Software Engineering \\
    \hline
    \rowcolor{rowaltbg}
    \textbf{Subject Code}      & CS3004 \\
   
    \hline
    \rowcolor{rowaltbg}
    \textbf{Date}              & \today \\
    \hline
    \rowcolor{rowaltbg}
    \textbf{Group Members}     &
      \begin{tabular}[t]{@{}L{4cm} L{3.5cm}@{}}
        Suyash Pandey        & : 2023UG1100 \\[3pt]
        Suraj Kumar          & : 2023UG1106 \\[3pt]
        Himanshu Raj          & : 2023UG1092 \\[3pt]
        Satyendra Kumar         & : 2023UG1105 \\
      \end{tabular} \\
    \hline
    \rowcolor{rowaltbg}
    \textbf{Institution}       & Indian Institute of Information Technology Ranchi, Jharkhand \\
    \hline
  \end{tabular}
\end{center}
\vfill

\newpage

% ── TABLE OF CONTENTS ─────────────────────────────────────────────────────────
\thispagestyle{fancy}
\tableofcontents
\newpage

% ══════════════════════════════════════════════════════════════════════════════
\section{Introduction}

\subsection{Purpose}

This Software Requirements Specification (SRS) document describes the functional and
non-functional requirements for the \textbf{Hotel Automation Software}. The purpose of
this document is to provide a comprehensive, unambiguous description of the system's
behaviour to guide development, testing, and validation of the software.

This software will automate the day-to-day operational activities of a 5-star hotel,
including room management, guest reservation, catering services, billing, and loyalty
programme management. The SRS serves as a contractual agreement between stakeholders and
the development team.

\subsection{Scope}

The Hotel Automation Software is designed to manage the complete operational workflow of a
5-star hotel, from room allocation and guest check-in through to catering, billing, and
check-out. The software will:

\begin{itemize}[leftmargin=1.5em]
  \item Maintain a database of all hotel rooms categorised by bed type (single / double)
        and amenity type (AC / Non-AC), along with their associated tariffs.
  \item Allow dynamic revision of room tariffs based on the prevailing occupancy rate,
        adjusted by a configurable percentage either upward or downward.
  \item Support advance and walk-in room reservations, subject to real-time room availability.
  \item Assign unique room numbers and token numbers to accommodated guests and generate
        apology messages when accommodation is unavailable.
  \item Record catering orders (food item, quantity, guest token, date and time) entered
        by the catering services manager.
  \item Generate itemised bills at check-out, including room charges and catering charges,
        and print the net balance payable.
  \item Maintain a frequent-guest identity database and apply applicable discounts to bills.
\end{itemize}

\noindent The software is \textbf{NOT} intended to:

\begin{itemize}[leftmargin=1.5em]
  \item Control physical room infrastructure (lighting, HVAC, locks) in real time.
  \item Handle payroll or HR management for hotel staff.
  \item Provide travel or tourism booking services beyond hotel accommodation.
\end{itemize}

\subsection{Definitions, Acronyms, and Abbreviations}

\begin{longtable}{|L{3.5cm}|L{10cm}|}
  \hline
  \thead{Term / Acronym} & \thead{Definition} \\
  \hline
  \endhead
  \hline
  \endfoot
  \rowcolor{rowaltbg}
  \textbf{SRS}                      & Software Requirements Specification \\
  \hline
  \textbf{AC}                       & Air-Conditioned --- denotes a room equipped with air conditioning \\
  \hline
  \rowcolor{rowaltbg}
  \textbf{Non-AC}                   & Non-Air-Conditioned --- denotes a room without air conditioning \\
  \hline
  \textbf{Room Tariff}              & The per-night charge for a specific room, which may vary by room type (single/double, AC/Non-AC) and by season / occupancy rate \\
  \hline
  \rowcolor{rowaltbg}
  \textbf{Occupancy Rate}           & The percentage of hotel rooms that are occupied at a given point in time; used by management to decide tariff adjustments \\
  \hline
  \textbf{Token Number}             & A unique identifier assigned to each guest at check-in; used to associate catering orders and other services with the correct guest \\
  \hline
  \rowcolor{rowaltbg}
  \textbf{Advance Reservation}      & A room booking made prior to the guest's arrival date \\
  \hline
  \textbf{Walk-in Reservation}      & A room booking made by a guest who arrives at the hotel without a prior booking \\
  \hline
  \rowcolor{rowaltbg}
  \textbf{Receptionist}             & Hotel staff member responsible for recording guest data, assigning rooms, and issuing token numbers \\
  \hline
  \textbf{Catering Services Manager} & Hotel staff member responsible for recording food and beverage orders against guest token numbers \\
  \hline
  \rowcolor{rowaltbg}
  \textbf{Frequent Guest}           & A returning guest who has been issued a loyalty identity number and is eligible for special discounts \\
  \hline
  \textbf{Identity Number}          & A unique loyalty identifier issued to frequent guests that entitles them to discount benefits \\
  \hline
  \rowcolor{rowaltbg}
  \textbf{GUI}                      & Graphical User Interface \\
  \hline
  \textbf{RDBMS}                    & Relational Database Management System --- used to persist room, guest, catering, and billing data \\
  \hline
  \rowcolor{rowaltbg}
  \textbf{SA/SD}                    & Structured Analysis / Structured Design \\
  \hline
  \textbf{UML}                      & Unified Modelling Language \\
  \hline
\end{longtable}

\subsection{Overview}

The remainder of this document is organised as follows:

\begin{itemize}[leftmargin=1.5em]
  \item \textbf{Section 2} provides the overall description of the system, including product
        perspective, functions, users, constraints, and assumptions.
  \item \textbf{Section 3} details the specific functional requirements of the system.
  \item \textbf{Section 4} details the non-functional requirements including performance,
        usability, reliability, and portability.
  \item \textbf{Section 5} contains the appendix with a use-case summary, a data
        dictionary, a traceability matrix, and the revision history.
\end{itemize}

\newpage

% ══════════════════════════════════════════════════════════════════════════════
\section{Overall Description}

\subsection{Product Perspective}

The Hotel Automation Software is a standalone desktop application implemented in
\textbf{Java} with a \textbf{Java Swing}-based graphical user interface, backed by a
relational database for persistent storage of rooms, guests, catering records, and
billing data. It does not interface with any external reservation portal or payment
gateway in this version.

The system interacts with the following entities:

\begin{itemize}[leftmargin=1.5em]
  \item \textbf{Hotel Manager:} Configures room types, tariffs, and discount policies;
        reviews occupancy reports and adjusts tariffs.
  \item \textbf{Receptionist:} Records guest arrival data, processes reservations,
        assigns rooms and token numbers, and initiates check-out billing.
  \item \textbf{Catering Services Manager:} Enters food and beverage orders against guest
        token numbers during the guest's stay.
  \item \textbf{Guest:} The end beneficiary; interacts with the system indirectly through
        the receptionist and catering staff.
\end{itemize}

\subsection{Product Functions}

At a high level, the system shall:

\begin{enumerate}[leftmargin=1.8em]
  \item Maintain a configurable inventory of rooms classified by bed type and AC type,
        each with an associated tariff.
  \item Display the average occupancy rate for any given month to support management
        tariff decisions.
  \item Allow the manager to revise room tariffs upward or downward by a specified
        percentage.
  \item Accept guest reservation requests (advance or walk-in) with arrival date, advance
        payment, approximate duration of stay, and room type preference.
  \item Allocate an available room and assign a unique token number to each accommodated
        guest; generate an apology message when no suitable room is available.
  \item Record catering orders (item, quantity, token number, date and time) entered by
        the catering services manager.
  \item Generate a complete itemised bill at check-out comprising room charges and
        catering charges, and display the net balance payable after deducting any advance
        payment.
  \item Issue loyalty identity numbers to frequent guests and apply the applicable
        discount to their bills.
\end{enumerate}

\subsection{User Classes and Characteristics}

\begin{longtable}{|L{3cm}|L{6cm}|L{4.5cm}|}
  \hline
  \thead{User Class} & \thead{Description} & \thead{Privilege Level} \\
  \hline
  \endhead
  \hline
  \endfoot
  \rowcolor{rowaltbg}
  \textbf{Hotel Manager}
    & Configures the system (rooms, tariffs, discounts) and reviews occupancy and revenue reports.
    & Administrative access \\
  \hline
  \textbf{Receptionist}
    & Handles guest reservations, room assignments, token issuance, and check-out processing.
    & Operational access \\
  \hline
  \rowcolor{rowaltbg}
  \textbf{Catering Services Manager}
    & Records food and beverage orders against guest token numbers.
    & Catering module access \\
  \hline
  \textbf{System Administrator}
    & Installs and maintains the application and database on hotel workstations.
    & Configuration and maintenance access \\
  \hline
\end{longtable}

\subsection{Operating Environment}

\begin{itemize}[leftmargin=1.5em]
  \item \textbf{Platform:} Desktop application (Windows 10+ / Ubuntu 20.04+ / macOS 11+)
  \item \textbf{Implementation Language:} Java (JDK 8 or above)
  \item \textbf{User Interface:} Java Swing
  \item \textbf{Database:} Embedded relational database (e.g., SQLite or Apache Derby)
  \item \textbf{Minimum RAM:} 512 MB
  \item \textbf{Minimum Disk Space:} 200 MB (including database storage)
  \item \textbf{Display:} Minimum 1024 $\times$ 768 resolution
  \item \textbf{Reference Hardware for Performance Targets:} Intel Core i5 (8th Gen or equivalent), 8 GB RAM, SSD storage
\end{itemize}

\subsection{Design and Implementation Constraints}

\begin{itemize}[leftmargin=1.5em]
  \item The system must be implemented entirely in Java.
  \item The user interface must be developed using Java Swing.
  \item All persistent data (rooms, guests, catering records, bills) must be stored in a
        relational database; no flat-file storage shall be used for transactional data.
  \item The application shall be delivered as a single executable JAR file bundling the
        embedded database driver.
  \item No external payment gateway or network connectivity is required in this version.
  \item Token numbers must be unique across the lifetime of the system, not merely within
        a single day.
\end{itemize}

\subsection{Assumptions and Dependencies}

\begin{itemize}[leftmargin=1.5em]
  \item Room tariffs are defined per room type (bed type $\times$ AC type) and are uniform
        within a type; individual room-level pricing is not supported in this version.
  \item Tariff adjustments apply to all rooms of the targeted type simultaneously.
  \item A guest is uniquely identified by their token number for the duration of their stay;
        the same guest may receive a new token number on each visit unless recognised as a
        frequent guest.
  \item Catering charges are computed from the quantities and unit prices stored in the
        system at the time of bill generation.
  \item The system assumes inputs are provided in good faith; basic input validation is
        required but adversarial inputs are not a primary concern.
  \item The ``approximate duration of stay'' entered at check-in is for planning purposes
        only; actual charges are computed from the real check-out date and time.
\end{itemize}

\newpage

% ══════════════════════════════════════════════════════════════════════════════
\section{Specific Functional Requirements}

\subsection{Room Management Requirements}

\subsubsection{Room Inventory Configuration}

\req{FR-1.1}{The system shall allow the hotel manager to define and maintain a list of
rooms, each characterised by a unique room number, bed type (single or double), and
AC type (AC or Non-AC).}

\req{FR-1.2}{The system shall associate a base tariff (per night, in currency units) with
each combination of bed type and AC type (i.e., up to four tariff categories: single-AC,
single-Non-AC, double-AC, double-Non-AC).}

\req{FR-1.3}{The system shall track the current availability status of each room
(\textit{available} or \textit{occupied}).}

\subsubsection{Occupancy Reporting and Tariff Revision}

\req{FR-2.1}{The system shall compute and display the average occupancy rate for any
given month, defined as:}
\[
  \text{Average Occupancy Rate (\%)} =
  \frac{\text{Total Room-Nights Occupied}}{\text{Total Room-Nights Available}} \times 100
\]

\req{FR-2.2}{The system shall allow the hotel manager to revise the tariff for any room
type by specifying a percentage adjustment (positive for upward revision, negative for
downward revision). The new tariff shall be computed as:}
\[
  \text{New Tariff} = \text{Current Tariff} \times \left(1 + \frac{\text{Adjustment \%}}{100}\right)
\]

\req{FR-2.3}{The system shall store a history of tariff changes (room type, old tariff,
new tariff, date of change) to support audit and reporting.}

% ─────────────────────────────────────────────────────────────────────────────
\subsection{Reservation Requirements}

\subsubsection{Guest Data Entry}

\req{FR-3.1}{The system shall allow the receptionist to enter the following data for each
reservation request:}
\begin{itemize}[leftmargin=1.5em]
  \item Guest name
  \item Arrival date and time
  \item Approximate duration of stay (in nights)
  \item Room type required (bed type and AC type)
  \item Advance payment amount (zero for walk-in guests with no advance)
  \item Frequent-guest identity number (optional; leave blank if not a frequent guest)
\end{itemize}

\req{FR-3.2}{The system shall support both advance reservations (made before the arrival
date) and walk-in reservations (made on the day of arrival).}

\subsubsection{Room Allocation}

\req{FR-4.1}{Upon receiving a reservation request, the system shall check the availability
of rooms matching the requested type on the requested dates.}

\req{FR-4.2}{If a suitable room is available, the system shall:}
\begin{itemize}[leftmargin=1.5em]
  \item Allocate the room with the lowest available room number of the requested type.
  \item Mark the room as \textit{occupied} for the duration of the stay.
  \item Assign a unique token number to the guest.
  \item Display and print the allocated room number and token number to the receptionist.
\end{itemize}

\req{FR-4.3}{If no suitable room is available, the system shall generate and display an
apology message to the receptionist indicating the unavailability and the reason (no
room of the requested type available on the requested dates).}

\req{FR-4.4}{The system shall not double-allocate a room to two overlapping reservations.}

% ─────────────────────────────────────────────────────────────────────────────
\subsection{Catering Services Requirements}

\req{FR-5.1}{The system shall allow the catering services manager to enter catering orders
comprising:}
\begin{itemize}[leftmargin=1.5em]
  \item Guest token number
  \item Food item name
  \item Quantity consumed
  \item Date and time of consumption
\end{itemize}

\req{FR-5.2}{The system shall validate that the token number entered corresponds to a
currently checked-in guest, and shall display an error message if not.}

\req{FR-5.3}{The system shall maintain a menu of food items with associated unit prices,
configurable by the hotel manager.}

\req{FR-5.4}{The system shall accumulate all catering orders for a guest throughout their
stay and associate them with the guest's token number for billing purposes.}

% ─────────────────────────────────────────────────────────────────────────────
\subsection{Billing and Check-Out Requirements}

\req{FR-6.1}{When a guest prepares to check out, the receptionist shall initiate the
check-out process by entering the guest's token number.}

\req{FR-6.2}{The system shall compute the total room charge as:}
\[
  \text{Room Charge} = \text{Tariff per Night (at time of check-in)} \times \text{Actual Nights Stayed}
\]

\req{FR-6.3}{The system shall compute the total catering charge as the sum of
(quantity $\times$ unit price) across all catering orders associated with the guest's
token number.}

\req{FR-6.4}{The system shall generate and display an itemised bill comprising:}
\begin{itemize}[leftmargin=1.5em]
  \item Guest name, room number, token number, check-in and check-out dates
  \item Itemised room charges (tariff per night $\times$ nights)
  \item Itemised catering charges (food item, quantity, unit price, subtotal per item)
  \item Total bill amount (room + catering)
  \item Advance payment deducted
  \item Frequent-guest discount deducted (if applicable)
  \item Net balance payable
\end{itemize}

\req{FR-6.5}{The system shall print the bill to the system's default printer upon
receptionist confirmation.}

\req{FR-6.6}{Upon successful check-out, the system shall mark the room as
\textit{available} from the check-out date.}

% ─────────────────────────────────────────────────────────────────────────────
\subsection{Frequent-Guest Loyalty Requirements}

\req{FR-7.1}{The system shall maintain a frequent-guest registry storing the guest's name,
contact information, and a unique system-generated identity number.}

\req{FR-7.2}{The hotel manager shall be able to define the discount percentage applicable
to frequent guests. The discount shall apply to the total bill (room + catering charges)
before deducting advance payments.}

\req{FR-7.3}{When a reservation is made using a valid frequent-guest identity number, the
system shall automatically apply the configured discount at check-out.}

\req{FR-7.4}{The system shall allow the receptionist to enrol a new frequent guest during
the check-out process, issuing a new identity number that is effective from the next stay.}

% ─────────────────────────────────────────────────────────────────────────────
\subsection{User Interface Requirements}

\req{FR-8.1}{The system shall provide a graphical user interface built using Java Swing,
with separate functional panels for Room Management, Reservations, Catering, Billing,
and Reports.}

\req{FR-8.2}{The GUI shall provide input forms for all reservation parameters described
in FR-3.1, with field-level validation and descriptive error messages.}

\req{FR-8.3}{The GUI shall provide a \textbf{Check Availability} button that queries room
availability before committing a reservation.}

\req{FR-8.4}{The GUI shall provide a \textbf{Confirm Reservation} button to finalise the
allocation and issue the token number.}

\req{FR-8.5}{The GUI shall provide a \textbf{Check-Out} button that initiates bill
generation for the guest identified by the entered token number.}

\req{FR-8.6}{The GUI shall provide a \textbf{Print Bill} button that sends the generated
bill to the system's default printer.}

\req{FR-8.7}{The GUI shall provide an \textbf{Adjust Tariff} panel accessible to the hotel
manager, allowing percentage-based tariff revisions per room type.}

\req{FR-8.8}{The GUI shall provide an \textbf{Occupancy Report} panel displaying the
average occupancy rate for a user-selected month.}

\req{FR-8.9}{The system shall validate all user inputs before processing and display
specific, actionable error messages for each invalid input field.}

\req{FR-8.10}{All input fields shall carry descriptive labels and tooltips explaining
expected values, formats, and valid ranges.}

\newpage

% ══════════════════════════════════════════════════════════════════════════════
\section{Non-Functional Requirements}

\subsection{Performance Requirements}

\begin{longtable}{|L{1.8cm}|L{3cm}|L{5.5cm}|L{3cm}|}
  \hline
  \thead{ID} & \thead{Requirement} & \thead{Description} & \thead{Metric} \\
  \hline
  \endhead
  \hline
  \endfoot
  \rowcolor{rowaltbg}
  \textbf{NFR-P1} & Query Response
    & Room availability queries and occupancy reports shall return results promptly on the
      reference hardware (Intel Core i5 8th Gen, 8 GB RAM) for a hotel with up to 500 rooms
      and 5 years of historical data.
    & $<$ 2 seconds \\
  \hline
  \textbf{NFR-P2} & Bill Generation
    & Generating and displaying the itemised check-out bill shall complete quickly,
      regardless of the number of catering line items.
    & $<$ 3 seconds \\
  \hline
  \rowcolor{rowaltbg}
  \textbf{NFR-P3} & UI Responsiveness
    & The GUI shall remain responsive during all database operations (operations execute
      on a background thread).
    & No EDT freeze $>$ 500\,ms \\
  \hline
  \textbf{NFR-P4} & Startup Time
    & The application shall start and be ready for user input promptly after launch,
      including database initialisation.
    & $<$ 5 seconds \\
  \hline
\end{longtable}

\subsection{Usability Requirements}

\begin{itemize}[leftmargin=1.5em]
  \item \textbf{NFR-U1:} The system shall provide a user interface that requires no
        programming knowledge to operate.
  \item \textbf{NFR-U2:} All input fields shall have descriptive labels and tooltips
        explaining expected values, formats, and valid ranges (see also FR-8.10).
  \item \textbf{NFR-U3:} Error messages shall be specific, identifying the offending
        field and the corrective action required.
  \item \textbf{NFR-U4:} Billing output shall be presented in a clearly formatted,
        itemised layout readable without domain expertise.
  \item \textbf{NFR-U5:} A first-time receptionist with basic computer skills shall be
        able to complete a reservation and check-out within 10 minutes without consulting
        documentation.
\end{itemize}

\subsection{Reliability Requirements}

\begin{itemize}[leftmargin=1.5em]
  \item \textbf{NFR-R1:} The system shall not crash or produce erroneous results for any
        valid input combination.
  \item \textbf{NFR-R2:} The system shall handle all boundary conditions gracefully,
        including zero available rooms, a guest with no catering orders, and a stay of
        zero nights.
  \item \textbf{NFR-R3:} No two guests shall be allocated the same room for any
        overlapping date range.
  \item \textbf{NFR-R4:} All financial calculations (tariff application, discount, net
        balance) shall be accurate to two decimal places.
  \item \textbf{NFR-R5:} The system shall preserve all reservation and billing records
        across application restarts; no transactional data shall be lost due to an
        unexpected shutdown after the transaction has been committed to the database.
\end{itemize}

\subsection{Maintainability Requirements}

\begin{itemize}[leftmargin=1.5em]
  \item \textbf{NFR-M1:} The source code shall follow standard Java coding conventions
        (Oracle Code Conventions) and every public class and method shall have a Javadoc
        comment.
  \item \textbf{NFR-M2:} The codebase shall be structured in a modular manner with clear
        separation between: database access layer, business logic, and UI components.
  \item \textbf{NFR-M3:} Adding a new room type or a new food item shall require only
        data-level changes (inserts into the room-type or menu table) with no modification
        to business logic code.
\end{itemize}

\subsection{Portability Requirements}

\begin{itemize}[leftmargin=1.5em]
  \item \textbf{NFR-PO1:} The software shall run on any platform supporting Java Runtime
        Environment (JRE 8 or above), including Windows 10+, Ubuntu 20.04+, and macOS 11+.
  \item \textbf{NFR-PO2:} The application shall be packaged as a single executable JAR
        file bundling the embedded database driver, requiring no external installation.
  \item \textbf{NFR-PO3:} The software shall not depend on any operating-system-specific
        libraries or absolute file paths.
\end{itemize}

\subsection{Security Requirements}

\begin{itemize}[leftmargin=1.5em]
  \item \textbf{NFR-S1:} The application shall support role-based access control, ensuring
        that catering staff cannot access billing or tariff management functions, and
        receptionists cannot access system configuration.
  \item \textbf{NFR-S2:} Exported or printed bills shall contain only guest transaction
        data and shall not include any system-level information (OS version, database
        paths, internal IDs).
  \item \textbf{NFR-S3:} Input validation (FR-8.9) shall prevent the application from
        entering an undefined or erroneous state due to malformed, empty, or out-of-range
        inputs.
  \item \textbf{NFR-S4:} Guest personal data (name, contact details) shall be stored only
        in the local database and shall not be transmitted to any external system.
\end{itemize}

\newpage

% ══════════════════════════════════════════════════════════════════════════════
\section{Appendix}

% ─────────────────────────────────────────────────────────────────────────────
\subsection{Use-Case Summary}
\label{sec:usecases}

\begin{longtable}{|L{1.5cm}|L{3cm}|L{3cm}|L{6cm}|}
  \hline
  \thead{UC ID} & \thead{Use Case} & \thead{Actor(s)} & \thead{Brief Description} \\
  \hline
  \endhead
  \hline
  \endfoot
  \rowcolor{rowaltbg}
  UC-01 & Make Reservation       & Receptionist        & Guest requests a room; receptionist enters details; system allocates room and issues token. \\
  \hline
  UC-02 & Handle No Availability & Receptionist        & No suitable room found; system generates apology message. \\
  \hline
  \rowcolor{rowaltbg}
  UC-03 & Record Catering Order  & Catering Mgr        & Catering manager enters food items against a guest token number. \\
  \hline
  UC-04 & Process Check-Out      & Receptionist        & Receptionist initiates check-out; system generates and prints itemised bill. \\
  \hline
  \rowcolor{rowaltbg}
  UC-05 & Adjust Room Tariff     & Hotel Manager       & Manager reviews occupancy report and revises tariff by a percentage. \\
  \hline
  UC-06 & View Occupancy Report  & Hotel Manager       & Manager selects a month; system displays average occupancy rate. \\
  \hline
  \rowcolor{rowaltbg}
  UC-07 & Enrol Frequent Guest   & Receptionist        & Receptionist enrols a qualifying guest; system issues a loyalty identity number. \\
  \hline
  UC-08 & Apply Loyalty Discount & Receptionist        & Frequent guest's identity number is entered; system applies configured discount at check-out. \\
  \hline
\end{longtable}

% ─────────────────────────────────────────────────────────────────────────────
\subsection{Data Dictionary}
\label{sec:datadictionary}

\begin{longtable}{|L{3.5cm}|L{2.5cm}|L{2cm}|L{5.5cm}|}
  \hline
  \thead{Attribute} & \thead{Data Type} & \thead{Constraints} & \thead{Description} \\
  \hline
  \endhead
  \hline
  \endfoot
  \rowcolor{rowaltbg}
  \texttt{room\_number}       & Integer    & PK, $\geq 1$       & Unique identifier for each physical room \\
  \hline
  \texttt{bed\_type}          & Enum       & \{Single, Double\} & Type of bed in the room \\
  \hline
  \rowcolor{rowaltbg}
  \texttt{ac\_type}           & Enum       & \{AC, Non-AC\}     & Air-conditioning category of the room \\
  \hline
  \texttt{tariff\_per\_night} & Decimal    & $> 0$              & Current nightly rate for the room type \\
  \hline
  \rowcolor{rowaltbg}
  \texttt{token\_number}      & Integer    & PK, unique, auto   & System-generated unique guest token \\
  \hline
  \texttt{guest\_name}        & String     & Non-empty          & Full name of the guest \\
  \hline
  \rowcolor{rowaltbg}
  \texttt{arrival\_datetime}  & DateTime   & Not null           & Check-in date and time \\
  \hline
  \texttt{checkout\_datetime} & DateTime   & $\geq$ arrival     & Actual check-out date and time \\
  \hline
  \rowcolor{rowaltbg}
  \texttt{advance\_paid}      & Decimal    & $\geq 0$           & Amount paid in advance by the guest \\
  \hline
  \texttt{identity\_number}   & Integer    & FK (optional)      & Frequent-guest loyalty ID; null if not enrolled \\
  \hline
  \rowcolor{rowaltbg}
  \texttt{food\_item}         & String     & Non-empty          & Name of the food or beverage item ordered \\
  \hline
  \texttt{quantity}           & Integer    & $\geq 1$           & Number of units of the food item consumed \\
  \hline
  \rowcolor{rowaltbg}
  \texttt{unit\_price}        & Decimal    & $> 0$              & Price per unit of the food item \\
  \hline
  \texttt{order\_datetime}    & DateTime   & Not null           & Date and time the catering order was placed \\
  \hline
  \rowcolor{rowaltbg}
  \texttt{discount\_pct}      & Decimal    & $0 \leq d \leq 100$& Percentage discount applied to frequent guests \\
  \hline
\end{longtable}

% ─────────────────────────────────────────────────────────────────────────────
\subsection{Requirements Traceability Matrix}

\begin{longtable}{|L{2.5cm}|L{5.5cm}|L{1.8cm}|L{1.8cm}|L{1.8cm}|}
  \hline
  \thead{Req.\ ID} & \thead{Requirement Description} & \thead{Priority} & \thead{Module} & \thead{Status} \\
  \hline
  \endhead
  \hline
  \endfoot
  \rowcolor{rowaltbg}
  FR-1.1 -- FR-1.3 & Room inventory, type classification, availability tracking
    & \textcolor{hred}{\textbf{High}} & Room Mgmt & Pending \\
  \hline
  FR-2.1 -- FR-2.3 & Occupancy reporting, tariff revision, change history
    & \textcolor{hred}{\textbf{High}} & Room Mgmt & Pending \\
  \hline
  \rowcolor{rowaltbg}
  FR-3.1 -- FR-3.2 & Guest data entry, advance and walk-in reservation support
    & \textcolor{hred}{\textbf{High}} & Reservation & Pending \\
  \hline
  FR-4.1 -- FR-4.4 & Room availability check, allocation, token issuance, apology message
    & \textcolor{hred}{\textbf{High}} & Reservation & Pending \\
  \hline
  \rowcolor{rowaltbg}
  FR-5.1 -- FR-5.4 & Catering order entry, validation, menu management, accumulation
    & \textcolor{hred}{\textbf{High}} & Catering & Pending \\
  \hline
  FR-6.1 -- FR-6.6 & Check-out initiation, room/catering charge computation, itemised bill, print, room release
    & \textcolor{hred}{\textbf{High}} & Billing & Pending \\
  \hline
  \rowcolor{rowaltbg}
  FR-7.1 -- FR-7.4 & Frequent-guest registry, discount config, discount application, enrolment
    & \textcolor{mred}{\textbf{Medium}} & Loyalty & Pending \\
  \hline
  FR-8.1 -- FR-8.10 & Java Swing GUI, reservation/catering/billing panels, validation
    & \textcolor{hred}{\textbf{High}} & UI Layer & Pending \\
  \hline
  \rowcolor{rowaltbg}
  NFR-R3 & No double-allocation of rooms for overlapping dates
    & \textcolor{hred}{\textbf{High}} & Reservation & Pending \\
  \hline
  NFR-R4 & Financial calculations accurate to two decimal places
    & \textcolor{hred}{\textbf{High}} & Billing & Pending \\
  \hline
  \rowcolor{rowaltbg}
  NFR-S1 & Role-based access control for staff user classes
    & \textcolor{mred}{\textbf{Medium}} & UI Layer & Pending \\
  \hline
\end{longtable}

% ─────────────────────────────────────────────────────────────────────────────
\subsection{Revision History}

\begin{longtable}{|L{2cm}|L{3cm}|L{4cm}|L{4.5cm}|}
  \hline
  \thead{Version} & \thead{Date} & \thead{Author} & \thead{Description} \\
  \hline
  \endhead
  \hline
  \endfoot
  \rowcolor{rowaltbg}
  0.1 & April 2026 & Project Team & Initial Draft \\
  \hline
  1.0 & \today     & Project Team & First complete draft; covers room management, reservation, catering, billing, and frequent-guest requirements; non-functional requirements; use-case summary; data dictionary; traceability matrix \\
  \hline
\end{longtable}

\vfill
\begin{center}
  \color{medblue}\rule{0.6\linewidth}{0.5pt}\\[0.5em]
  \color{gray}\small\textit{End of Document --- Hotel Automation Software SRS v1.0}
\end{center}

\end{document}